% 原書前付の全訳。原書PDFのページ番号を各項目に記す。
% 原書の目次（PDF 11--13）は本文の見出しから別途生成する。
% 原書の扉（PDF 7）の題名・著者・所属は親原稿のタイトルに収録する。

% 原書PDF 3。
\sectionlink*{本書について}
\begin{center}
\textbf{弦理論}\\
ボソン弦入門
\end{center}

『弦理論』を構成するこの2巻は、弦理論への最新かつ包括的な、
教育的配慮の行き届いた入門を提供する。

第I巻『ボソン弦入門』では、Polyakov経路積分と共形場理論に基づき、
ボソン弦を詳しく導入する。最初の4章では、弦理論の中心的な考え方、
共形場理論およびPolyakov経路積分の手法、弦の共変量子化を導入する。
続く3章では弦の相互作用を扱い、その一般的な形式を述べたうえで、
木レベル振幅と1ループ振幅を詳しく取り上げる。
第8章では、トーラスコンパクト化と、$T$双対性やDブレーンをはじめとする
弦の物理の重要な諸側面を扱う。第9章では、有限性とユニタリティの解析を含む
高次振幅、および非摂動的なさまざまな考え方を扱う。
経路積分法の短い講義を収めた付録も含まれている。

第II巻『超弦理論とその先へ』は、超対称弦理論の導入から始まり、
近年の重要な進展を幅広く紹介する。最初の3章では、
I型、II型、ヘテロティック超弦理論とそれらの相互作用を導入する。
続く2章では、強く結合した弦についての近年の重要な発見を取り上げる。
Dブレーンとそのダイナミクスの詳しい議論に始まり、
弦の双対性、M理論、ブラックホールのエントロピーを扱う。
その次の章では、共形場理論の古典的な成果を数多くまとめる。
最後の4章は4次元弦理論を扱い、二つの目的を持つ。
一つは、最も単純な弦模型のいくつかが標準模型を統一するための従来の考え方と
どのようにつながるかを示すことであり、もう一つは、
世界面と時空の対称性についての重要で美しい一般的な成果を数多くまとめることである。
付録では、フェルミオンと超対称性について必要となる予備知識を要約する。

両巻とも、学生に有用な文献を重視した注釈付きの参考文献一覧と、
重要な用語・概念についての詳細な用語集を収録している。
本文の要点を確実に身につけ、さらに新しい考え方を取り入れるための演習問題も
数多く用意されている。

現代の超弦理論に関心を持つ理論物理学、素粒子物理学、相対論の大学院生と研究者にとって、
必携の教科書であり参考書である。

Joseph Polchinskiは、1980年にUniversity of California at Berkeleyで
博士号を取得した。Stanford Linear Accelerator CenterとHarvardで
博士研究員を務めたのち、1984年にUniversity of Texas at Austinの
教員となり、1992年に現在の職であるUniversity of California at Santa Barbaraの
物理学教授、およびInstitute for Theoretical Physicsの常任研究員に就任した。

Polchinski教授は、明快で教育的な解説に秀でているだけでなく、
第一線の弦理論研究者でもある。1995年にDブレーンの重要性を発見したことは、
この分野における近年の最も重要な貢献の一つである。
また、場の量子論の多くの領域と、素粒子物理学の超対称模型にも重要な貢献をしている。

% 原書PDF 5。引用文の省略記号は原書にあるもの。
\clearpage
\sectionlink*{ハードカバー版の書評より}
\subsectionlink*{第1巻}
\begin{quote}
「……これは見事な本である。一貫した展開と、
各主題を扱う際の明快さと洞察の深さが際立っている。……
大学院の高度な分野でこれ以上によい教科書を思い浮かべることは難しく、
その分野の大家自身が書いた教科書に出会えるのはまれである。
本書には、有用な問題集、用語集、そして並外れて充実した索引も
含まれていることを指摘しておきたい。」
\begin{flushright}Physics Today\end{flushright}
\end{quote}

\begin{quote}
「……『現代的な』弦理論の始まりを画した、1990年代前半から半ばにかけての
超弦革命の発見を扱う、最も包括的な教科書である。」
\begin{flushright}
Donald Marolf, University of California, Santa Barbara\\
American Journal of Physics
\end{flushright}
\end{quote}

\begin{quote}
「物理学者は、量子力学と一般相対論および素粒子理論の統一を含む、
自然の根源的な理論への最も有望な道が、弦理論にあると考えている。
この優美な数理物理学の分野を、Joseph Polchinskiが
Cambridge University Press刊の2巻で詳しく解説している。……
上級の学生と研究者のために書かれたこの2巻は、徹底した最新の知識を提供する。」
\begin{flushright}American Scientist\end{flushright}
\end{quote}

\begin{quote}
「本書の教育的価値を強調したい。採られているアプローチは現代的で、
心地よいほど体系的であり、広範な成果を統一された言葉で扱っている。
各章末の演習問題が本文の議論を補っている。
他方、超弦の文脈で不可欠な手法と概念を導入しているため、
この分野の研究者にとっても有用な参考書である。」
\begin{flushright}Mathematical Reviews\end{flushright}
\end{quote}

\begin{quote}
「弦理論を志す人々が長い道のりを歩むうえで励みとなる本への需要を十分に満たしており、
古典となる運命にある。これは実に胸躍る営みであり、
この見事な本はそれに大いに貢献している。」
\begin{flushright}David Bailin, The Times Higher Education Supplement\end{flushright}
\end{quote}

\subsectionlink*{第2巻}
\begin{quote}
「要するに、この2巻は……現代の超弦理論がもたらしうるさまざまな帰結に関心を持つ、
素粒子物理学と相対論の学生および研究者にとって、
標準的な教科書・参考書となるであろう。」
\begin{flushright}Allen C. Hirshfeld, General Relativity and Gravitation\end{flushright}
\end{quote}

\begin{quote}
「Polchinskiは、この4年間に弦理論の理解を一変させた刺激的な進展の主要な立役者であり、
またSteven Weinbergが序文で証言しているように、模範的な教師でもある。
彼は、各章に数多くの演習問題を添えた、傑出した2巻の教科書を生み出した。
古典となる運命にある。……見事である。」
\begin{flushright}David Bailin, The Times Higher Education Supplement\end{flushright}
\end{quote}

\begin{quote}
「本巻は、超弦のダイナミクスについての現在の知識を、
詳細でありながら包括的に解説することに成功している。
扱う主題は、理論の基本的構成から、その非摂動的振る舞いに関する最新の発見にまで及ぶ。
議論は驚くほど自己完結しており
（本巻には、さまざまな次元におけるスピノルと超対称性についての有用な付録まである）、
この分野への入門書としても、この分野の研究者にとっての優れた参考書としても役立つであろう。」
\begin{flushright}Mathematical Reviews\end{flushright}
\end{quote}

% 原書PDF 6。書名・人名は書誌表記として英語を保持。
\clearpage
\sectionlink*{Cambridge Monographs on Mathematical Physics}
\noindent 総編集者：P.~V.~Landshoff、D.~R.~Nelson、S.~Weinberg
\begingroup
\small
\begin{list}{}{\setlength{\leftmargin}{0pt}\setlength{\rightmargin}{0pt}
\setlength{\itemsep}{2pt}\setlength{\parsep}{0pt}}
\item S.~J.~Aarseth, \emph{Gravitational N-Body Simulations}
\item J.~Ambjørn, B.~Durhuus and T.~Jonsson, \emph{Quantum Geometry: A Statistical Field Theory Approach}
\item A.~M.~Anile, \emph{Relativistic Fluids and Magneto-Fluids}
\item J.~A.~de Azcárrage and J.~M.~Izquierdo, \emph{Lie Groups, Lie Algebras, Cohomology and Some Applications in Physics}\textsuperscript{\(\dagger\)}
\item O.~Babelon, D.~Bernard and M.~Talon, \emph{Introduction to Classical Integrable Systems}
\item F.~Bastianelli and P.~van Nieuwenhuizen, \emph{Path Integrals and Anomalies in Curved Space}
\item V.~Belinkski and E.~Verdaguer, \emph{Gravitational Solitons}
\item J.~Bernstein, \emph{Kinetic Theory in the Expanding Universe}
\item G.~F.~Bertsch and R.~A.~Broglia, \emph{Oscillations in Finite Quantum Systems}
\item N.~D.~Birrell and P.~C.~W.~Davies, \emph{Quantum Fields in Curved Space}\textsuperscript{\(\dagger\)}
\item M.~Burgess, \emph{Classical Covariant Fields}
\item S.~Carlip, \emph{Quantum Gravity in \(2+1\) Dimensions}
\item J.~C.~Collins, \emph{Renormalization}\textsuperscript{\(\dagger\)}
\item M.~Creutz, \emph{Quarks, Gluons and Lattices}\textsuperscript{\(\dagger\)}
\item P.~D.~D'Eath, \emph{Supersymmetric Quantum Cosmology}
\item F.~de Felice and C.~J.~S.~Clarke, \emph{Relativity on Curved Manifolds}\textsuperscript{\(\dagger\)}
\item B.~S.~DeWitt, \emph{Supermanifolds, 2nd edition}\textsuperscript{\(\dagger\)}
\item P.~G.~O.~Freund, \emph{Introduction to Supersymmetry}\textsuperscript{\(\dagger\)}
\item J.~Fuchs, \emph{Affine Lie Algebras and Quantum Groups}\textsuperscript{\(\dagger\)}
\item J.~Fuchs and C.~Schweigert, \emph{Symmetries, Lie Algebras and Representations: A Graduate Course for Physicists}\textsuperscript{\(\dagger\)}
\item Y.~Fujii and K.~Maeda, \emph{The Scalar--Tensor Theory of Gravitation}
\item A.~S.~Galperin, E.~A.~Ivanov, V.~I.~Orievetsky and E.~S.~Sokatchev, \emph{Harmonic Superspace}
\item R.~Gambini and J.~Pullin, \emph{Loops, Knots, Gauge Theories and Quantum Gravity}\textsuperscript{\(\dagger\)}
\item M.~Göckeler and T.~Schücker, \emph{Differential Geometry, Gauge Theories and Gravity}\textsuperscript{\(\dagger\)}
\item C.~Gómez, M.~Ruiz Altaba and G.~Sierra, \emph{Quantum Groups in Two-Dimensional Physics}
\item M.~B.~Green, J.~H.~Schwarz and E.~Witten, \emph{Superstring Theory, volume 1: Introduction}\textsuperscript{\(\dagger\)}
\item M.~B.~Green, J.~H.~Schwarz and E.~Witten, \emph{Superstring Theory, volume 2: Loop Amplitudes, Anomalies and Phenomenology}\textsuperscript{\(\dagger\)}
\item V.~N.~Gribov, \emph{The Theory of Complex Angular Momenta}
\item S.~W.~Hawking and G.~F.~R.~Ellis, \emph{The Large-Scale Structure of Space-Time}\textsuperscript{\(\dagger\)}
\item F.~Iachello and A.~Arima, \emph{The Interacting Boson Model}
\item F.~Iachello and P.~van Isacker, \emph{The Interacting Boson--Fermion Model}
\item C.~Itzykson and J.-M.~Drouffe, \emph{Statistical Field Theory, volume 1: From Brownian Motion to Renormalization and Lattice Gauge Theory}\textsuperscript{\(\dagger\)}
\item C.~Itzykson and J.-M.~Drouffe, \emph{Statistical Field Theory, volume 2: Strong Coupling, Monte Carlo Methods, Conformal Field Theory, and Random Systems}\textsuperscript{\(\dagger\)}
\item C.~Johnson, \emph{D-Branes}
\item J.~I.~Kapusta, \emph{Finite-Temperature Field Theory}\textsuperscript{\(\dagger\)}
\item V.~E.~Korepin, A.~G.~Izergin and N.~M.~Boguliubov, \emph{The Quantum Inverse Scattering Method and Correlation Functions}\textsuperscript{\(\dagger\)}
\item M.~Le Bellac, \emph{Thermal Field Theory}\textsuperscript{\(\dagger\)}
\item Y.~Makeenko, \emph{Methods of Contemporary Gauge Theory}
\item N.~Manton and P.~Sutcliffe, \emph{Topological Solitons}
\item N.~H.~March, \emph{Liquid Metals: Concepts and Theory}
\item I.~M.~Montvay and G.~Münster, \emph{Quantum Fields on a Lattice}\textsuperscript{\(\dagger\)}
\item L.~O' Raifeartaigh, \emph{Group Structure of Gauge Theories}\textsuperscript{\(\dagger\)}
\item T.~Ortín, \emph{Gravity and Strings}
\item A.~Ozorio de Almeida, \emph{Hamiltonian Systems: Chaos and Quantization}\textsuperscript{\(\dagger\)}
\item R.~Penrose and W.~Rindler, \emph{Spinors and Space-Time, volume 1: Two-Spinor Calculus and Relativistic Fields}\textsuperscript{\(\dagger\)}
\item R.~Penrose and W.~Rindler, \emph{Spinors and Space-Time, volume 2: Spinor and Twistor Methods in Space-Time Geometry}\textsuperscript{\(\dagger\)}
\item S.~Pokorski, \emph{Gauge Field Theories, 2nd edition}
\item J.~Polchinski, \emph{String Theory, volume 1: An Introduction to the Bosonic, String}\textsuperscript{\(\dagger\)}
\item J.~Polchinski, \emph{String Theory, volume 2: Superstring Theory and Beyond}\textsuperscript{\(\dagger\)}
\item V.~N.~Popov, \emph{Functional Integrals and Collective Excitations}\textsuperscript{\(\dagger\)}
\item R.~J.~Rivers, \emph{Path Integral Methods in Quantum Field Theory}\textsuperscript{\(\dagger\)}
\item R.~G.~Roberts, \emph{The Structure of the Proton}\textsuperscript{\(\dagger\)}
\item C.~Rovelli, \emph{Quantum Gravity}
\item W.~C.~Saslaw, \emph{Gravitational Physics of Stellar and Galactic Systems}\textsuperscript{\(\dagger\)}
\item H.~Stephani, D.~Kramer, M.~A.~H.~MacCallum, C.~Hoenselaers and E.~Herlt, \emph{Exact Solutions of Einstein's Field Equations, 2nd edition}
\item J.~M.~Stewart, \emph{Advanced General Relativity}\textsuperscript{\(\dagger\)}
\item A.~Vilenkin and E.~P.~S.~Shellard, \emph{Cosmic Strings and Other Topological Defects}\textsuperscript{\(\dagger\)}
\item R.~S.~Ward and R.~O.~Wells Jr, \emph{Twistor Geometry and Field Theories}\textsuperscript{\(\dagger\)}
\item J.~R.~Wilson and G.~J.~Mathews, \emph{Relativistic Numerical Hydrodynamics}
\end{list}
\noindent \textsuperscript{\(\dagger\)}ペーパーバック版刊行。
\endgroup

% 原書PDF 8。
\clearpage
\sectionlink*{刊行情報}
\begingroup
\small
\noindent\textbf{Cambridge University Press}\\
Cambridge, New York, Melbourne, Madrid, Cape Town, Singapore, São Paulo

\noindent Cambridge University Press\\
The Edinburgh Building, Cambridge CB2 8RU, UK\\
アメリカ合衆国ではCambridge University Press, New Yorkが刊行。\\
\href{https://www.cambridge.org}{\nolinkurl{www.cambridge.org}}\\
本書についての情報：\href{https://www.cambridge.org/9780521633031}{\nolinkurl{www.cambridge.org/9780521633031}}

\noindent\copyright\ Cambridge University Press 2001, 2005

本出版物は著作権により保護されている。法定の例外および関連する包括的利用許諾契約の
規定による場合を除き、Cambridge University Pressの書面による許可なく、
そのいかなる部分も複製することはできない。

\noindent 印刷版初版刊行：1998年

\begin{tabular}{@{}ll@{}}
ISBN-13 978-0-511-33821-2 & 電子書籍（NetLibrary）\\
ISBN-10 0-511-33821-X & 電子書籍（NetLibrary）\\[1ex]
ISBN-13 978-0-521-63303-1 & ハードカバー\\
ISBN-10 0-521-63303-6 & ハードカバー\\[1ex]
ISBN-13 978-0-521-67227-6 & ペーパーバック\\
ISBN-10 0-521-67227-9 & ペーパーバック
\end{tabular}

Cambridge University Pressは、本出版物中で言及される外部または第三者の
インターネット上のウェブサイトについて、そのURLの存続や正確さに責任を負わず、
そのようなウェブサイトの内容が正確または適切であること、
あるいは今後もそうあり続けることを保証しない。
\endgroup

% 原書PDF 9。
\bigskip
\begin{center}
\itshape 両親に捧ぐ
\end{center}

% 原書PDF 14。
\clearpage
\sectionlink*{第II巻の構成}
\begin{description}
\item[第10章] I型およびII型超弦
\item[第11章] ヘテロティック弦
\item[第12章] 超弦の相互作用
\item[第13章] Dブレーン
\item[第14章] 強結合での弦
\item[第15章] 共形場理論の発展的事項
\item[第16章] オービフォールド
\item[第17章] Calabi--Yauコンパクト化
\item[第18章] 4次元の物理
\item[第19章] 発展的な話題
\item[付録B] さまざまな次元におけるスピノルと超対称性
\end{description}

% 原書PDF 15。
\clearpage
\sectionlink*{序文}
素粒子の標準模型は、その成功にもかかわらず、
強い相互作用と電弱相互作用に加えて重力をも取り込む、
より広い理論の中に組み込まれなければならないことが、
当初から明らかであった。現在、そのような理論のもっともらしい候補は一つしかない。
それが弦の理論である。この理論は、1960年代にハドロンのあまり成功しなかった模型として
出発し、のちになってようやく、すべての力の理論となる可能性を持つものとして姿を現した。

読者を弦理論へ導くうえで、Joseph Polchinskiほど適任な人はいない。
その理由の一つは、彼がこの理論の発展に重要な役割を果たしてきたことである。
最近の例を一つだけ挙げると、彼は「Dirichletブレーン」という新しい種類の
拡がりを持つ物体の可能性を発見した。
これは、それまで異なるものと考えられていた弦理論どうしの関係を解き明かす、
ここ数年の目覚ましい進展に不可欠な要素となっている。

それと同じく重要なのは、Polchinskiが、複雑な数学的形式の中で
何が物理的に重要かを見抜き、それを他の人に説明するという、
まれに見る才能を備えていることである。
本書の校正刷りに目を通しながら、私はPolchinskiが
University of Texas at AustinのTheory Groupの一員であったころのことを思い出した。
弦理論の中で私を悩ませていた点について、
彼の忍耐強く明快な説明から何度も恩恵を受けたのである。
この刺激的な分野を修得したいと願う、すべての物理学者に本書を推薦する。

\begin{flushright}
Steven Weinberg\\
シリーズ編集者\\
Cambridge Monographs on Mathematical Physics\\
1998年
\end{flushright}

% 原書PDF 17--19。
\clearpage
\sectionlink*{まえがき}
10年以上前、私が弦理論の本を書こうと初めて決心したころは、
学生時代の記憶は今よりもずっと鮮明だった。
それでも、難しい主題を理解しやすくしてくれる教科書に出会うことが
何よりの喜びの一つだったことは覚えているし、
弦理論についてもそのような本を提供したいと思っていた。

したがって、私の第一の目的は、Polyakov経路積分と共形場理論に基づく、
首尾一貫した弦理論の入門書を提供することだった。
弦理論についての予備知識はいっさい仮定しない。
ただし、読者が計量や曲率などの一般相対論の中心的な考え方、
および非Abelゲージ対称性までの場の量子論の考え方に精通していることは仮定する。
当初は、場の量子論の一通りの講義を履修していることを前提としていたが、
できるだけ早く弦理論を学びたいと熱望する学生が多いこと、
また、場の量子論を履修していても、その講義では弦理論に必要な手法が
重視されていなかったという学生もいることが明らかになった。
そこで、そうした手法については自己完結した導入を与えるよう努めた。

第二の目的は、最も単純な4次元弦理論のいくつかが、
標準模型を統一するための従来の考え方とどのようにつながるかを示すこと、
そして、世界面と時空の対称性から導かれる4次元弦理論の物理に関する
一般的な成果をまとめることであった。
新たな進展によって第三の目標も加わった。
それは、弦の双対性、M理論、Dブレーン、ブラックホールのエントロピーに関する
最近の発見を紹介することである。

このような教科書を書くとき、完全な説明を与える必要と、
最も興味深い最近の成果にできるだけ早く到達したいという願いは衝突する。
私は、その両方を満たすよう努めた。
たとえば完全性の側面では、第6章のさまざまな経路積分を3通りの異なる方法で計算し、
ボソン弦の臨界次元を本文と演習問題を通じて7通りの異なる方法で計算している。
効率の側面では、この2巻をたどるいくつかの近道を以下に提案する。

特に問題となるのが弦の摂動論である。
この枠組みは必然的に第I巻の中心的な主題となるが、
最近の非摂動的な進展にとっては、やや副次的な位置にある。
そこでは、自由弦のスペクトルと時空の対称性のほうが、より決定的に重要である。
幸い、弦の摂動論からは、$T$双対性とDブレーンを経由して、
最近の発見へと至る自然な道筋がある。

可能な講義の一例は、第1--3章、第4.1節、第5--8章
（第5.4節と第6.7節を省く）、第10章、第11.1節、第11.2節、第11.6節、
第12.1節、第12.2節、そして第13章と第14章からなる。
私は、この順序なら2クォーターの授業で扱えると考えており、
弦理論の導入から、弦の双対性、M理論、最も単純なブラックホールのエントロピー計算まで
到達できる。さらに別の近道を第5.1節の末尾で提案している。

$T$双対性とそれに関連する弦特有の現象に関心のある読者は、
第4.1節から直接第8章へ進むことができる。
開弦の議論を追うためには、第6.5節の冒頭にあるChan--Paton因子の導入が必要であり、
Dブレーンの張力の計算を追うためには、
第7章で求める1ループ真空振幅が必要となる。

超対称弦に関心のある読者は、第4.1節のあとで第10章と第11章の大部分を読むことができる。
この場合も、開弦の議論を追うためにはChan--Paton因子の導入が必要であり、
第10.7節、第10.8節、第11.2節の整合性条件を追うためには1ループ真空振幅が必要となる。

共形場理論に関心のある読者は、第2章、第6.1節、第6.2節、第6.7節、
第7.1節、第7.2節、第8.2節、第8.3節（共形場理論の側面に注目する）、
第8.5節、第10.1--10.4節、第11.4節、第11.5節、そして第15章を読むとよいだろう。
4次元弦理論に関心のある読者は、第8章、第10章、第11章を読めば、
第16--19章の大部分を理解できる。

弦理論のように活発な分野では、ある見積もりでは文献数が1万編に近づいており、
ごく簡単にしか扱えない重要な主題も、まったく扱えない主題も必然的に生じる。
そのうちのいくつかについては、各巻末の参考文献一覧で総説論文を紹介している。
おそらく最も大きな省略は、曲がった多様体上でのコンパクト化を、
より完全には扱っていないことであろう。
この主題の幾何学的手法は、ここで重視する場の量子論の手法とはいささか方向が異なるため、
第17章と第19章では最も重要な成果の概要だけを収めた。
Green、Schwarz、Witten（1987）\cite{Green:1987mn}の第2巻には、もっと詳しい導入がある。
しかし、この主題はその後も重要性を増し続けており、
それだけを扱う入門書に明らかに値するものである。

本書は、1987--88年にUniversity of Texas at Austinで行った講義から生まれた。
当初の計画では、1年をかけて講義ノートを本にまとめるはずだった。
しかし、説明をより明快で完全なものにしたいという思いと、
研究に気を取られたことが、その妨げとなった。
初期の企画書では、完成予定日を1989年6月$\pm$1か月と見積もっていたが、
100標準偏差も外れてしまった。
8年の間、完成予定日はつねにおよそ1年先のままであり、
その間に1巻が2巻へと膨れ上がった。
幸いにも、ついに、その締切の一つが先延ばしにならずに済んだ。

私はまた、本書の一部をUniversity of California at Santa Barbaraでの講義、
および1994年のLes Houches、1995年のTrieste、1996年のTASIの各スクールでも用いた。
Nathan SeibergとMichael Douglas（Rutgers）、
Steven Weinberg（Texas）、Andrew StromingerとJuan Maldacena（Harvard）、
Nathan Berkovits（São Paulo）、Martin Einhorn（Michigan）も、
それぞれの講義で本書の一部を使用した。
これらの同僚とその学生たちから寄せられた、たいへん有用な意見に感謝したい。
また、この計画の初めに助言と励ましをくださったSteven Weinberg、
原稿全体を入念に読んでくださったShyamoli Chaudhuriにも感謝したい。
さらに、UT AustinとUC Santa BarbaraのDepartments of Physics、
UC Santa BarbaraのInstitute for Theoretical Physics、
National Science Foundationの支援に謝意を表したい。

本書の執筆が長期に及ぶ間に、何十人もの同僚が、
ここで扱う主題についての私の理解を明確にするのを助けてくださり、
何十人もの学生が訂正やその他の改善を提案してくれた。
両方のグループの人々の名を列挙しようとし始めたが、それが不可能なことに気づいた。
長くても不完全な一覧をここに掲げる代わりに、
正誤表のウェブサイト
\begin{center}
\url{http://www.itp.ucsb.edu/~joep/bigbook.html}
\end{center}
で、最新の一覧を維持することにする。

加えて、弦理論の発展に貢献してきたすべての人々に、まとめて感謝したい。
特に第II巻は、その大部分が、多くの物理学者によって導かれた美しい成果の集成であると
私には思われる。弦理論は、そしてその基礎となった物理学の全体もまた、
人類の偉大な達成の一つであり、
その現在の姿を捉えようと試みることができたのは、私にとって光栄であった。

最後に、これほどの規模の計画を完成させるには多くの犠牲が必要であり、
家族もそれを分かち合ってくれた。
Dorothy、Steven、Danielの理解、忍耐、支えに感謝したい。

\begin{flushright}
Joseph Polchinski\\
Santa Barbara, California\\
1998年
\end{flushright}

% 原書PDF 20--21。太字のS・Rと通常のS・Rの区別を画像で照合済み。
\clearpage
\sectionlink*{記法}
本書では、Misner、Thorne、Wheeler（1973）\cite{Misner:1973prb}の$+++$の規約を用いる。
特に、計量の符号は$(-++\cdots+)$である。
定数$\hbar$と$c$は1に置くが、Reggeスロープ$\alpha'$は明示しておく。

世界面座標およびモジュライ（$z$、$\tau$、$q$など）の共役を表すには上線
$\bar{\phantom{z}}$を用いるが、長い式には星印${}^{*}$を用いる。
時空のフェルミオン場につけた上線はDirac随伴を表し
（これは第II巻にのみ現れる）、世界面演算子につけた上線は
Euclid随伴（第6.7節で定義する）を表す。
弦上の自由度については、次の用語を同義として扱う。
\[
 \text{正則}=\text{左向きに伝播する},\qquad
 \text{反正則}=\text{右向きに伝播する}.
\]
これについては第2.1節で説明する。
本書の規約では、ヘテロティック弦の超対称な側は右向きに伝播する。
反正則演算子にはチルダ$\tilde{\phantom{z}}$をつける。
第2.3節で説明するように、これらは正則演算子の随伴ではない。
次の規約にも注意されたい。
\[
 d^{2}z\equiv 2\,dx\,dy,\qquad
 \delta^{2}(z,\bar z)\equiv\frac12\,\delta(x)\delta(y),
\]
ここで、$z=x+iy$は任意の複素変数である。
これらは、両方の定義の係数を1とする大部分の文献の規約とは異なる。

時空の作用は$\boldsymbol{S}$、世界面の作用は$S$と書く。
この区別はDブレーンについては問題を生む。
Dブレーンは前者と$T$双対であり、後者と$S$双対だからである。
ここでは任意に$\boldsymbol{S}$を用いることにした。
時空の計量は$G_{\mu\nu}$、世界面の計量は
$\gamma_{ab}$（Minkowski型）または$g_{ab}$（Euclid型）である。
第I巻では、時空のRicciテンソルは$\boldsymbol{R}_{\mu\nu}$、
世界面のRicciテンソルは$R_{ab}$とする。
第II巻では前者が頻繁に現れ、後者はまったく現れないため、
時空のRicciテンソルを$R_{\mu\nu}$と書くように変更している。

以下の記号を用いる。
\begin{description}
\item[$\equiv$] 左辺を右辺によって定義する。
\item[$\cong$] 同値である。
\item[$\approx$] 近似的に等しい。
\item[$\sim$] （演算子積展開では）非特異な項を除いて等しいこと、
または大まかな対応を表す。
\end{description}
